% !TeX root = ../papers/nc-p3.tex
\section{直線と第二曲線に制限した反対切断の余次元}
\label{sec:branch-rank-one-residual}
\begin{theorem}\label{thm:branch-rank-one-residual}
曲面射影のない標準正則窓の順中心 $p$ の中央支持次元を $a$、その後続直線 $L$ の
一般点の中央支持次元を $h$ とする。$h=4$ なら $\rho\ge a+6$ である。
二直線列では両中心の中央支持次元 $a,b$ と $r=\operatorname{rank} w(p,z)$ が
$r<\min(a,b)$ を満たす。特に反対階数8で二直線列は存在しない。
全て逆向きにも成り立つ。
\end{theorem}
\begin{proof}
二直線列の第二直線を $C$ とする。命題\ref{prop:two-line-branch-block}の
直交束の列によって $R_2$ の $C$ への制限は単射である。
従って $w(p)$ の中央支持は $C$ の消去空間に含まれない。
実際そうなら $w(p)$ の位置1からの縮約は全てこの単射の核に入る。
$w(p,z)$ の位置2支持は $w(p)$ の中央支持と $C$ の消去空間の交叉に入るので
$r<a$。逆向きに $r<b$。命題の $\rho\ge6+a+b-r$ と $a,b\ge2$ から下界9を得る。

一般の分岐列に戻り、$h=4$, $\rho=a+5$ と仮定する。これより小さい値は
分岐列の隣接中央支持の狭義不等式で除かれる。
第二曲線の次数は $d=3-f_1\in\{1,2,3\}$ である。
二項変数 $s,t$ を独立にして直線 $x(s)$ と第二曲線 $y(t)$ に制限すると
\[
             w(-,x(s),y(t),-)=[s,t]A(t)
\]
であり、$A$ は次数 $d-1$ の位置0と3の間の行列である。対角上の消滅と
$s$ に関する一次性を用いた。各一般 $t$ で $w(-,y(t))$ の全行は
$W_{x(t)}^\perp$ に含まれ、その $W_p$ による商は一次元である。
$W_p$ は全 $x(s)$ 上で消える。従って $s\ne t$ として評価すると
$A(t)$ の階数は高々1である。$d=1$ なら次数式 $d_3=-f_2$ より動く位置は二つだけとなるが、
二直線列の定数 $A$ は命題\ref{prop:two-line-branch-block}により階数3である。

$d=2$ の場合、第二曲線は平面 $\mathbb P(H)$ 内の滑らかな二次曲線である。
$H\subset V_2^*$ の次元は3で、二次の二項式の完全系によるパラメータを取る。
$w(x(s))=\sum_{j=1}^3 K_j(s)\otimes B_j$ とし、一次列 $K_j(s)$ を $H$ 上に制限する。
これらは完全二次系の一次関係であり、その二空間は
$(-t,s,0),(0,-t,s)$ 型の標準二列で張られ、各点で独立である。
三つの制限はこの二空間を張る。実際元の三列は一般点で独立なので、
一次元核 $H^\perp$ を法としても一般階数は2以上である。
そこで定数の列変換により最初の二列が標準二列に制限され、第三列は
$H^\perp$ の固定ベクトルと非零一次式の積となる。その一次式の零点 $q$ では
$K_j(q)$ の包から $H^*$ への射は単射である。
ゆえに $w(q)$ の位置0,3の支持は、その位置2を $H$ に制限しても保たれる。
上の表示の $A(t)$ は階数1以下の一次鉛筆なので、固定した左または右の一次元支持を持つ。
これは二つの可分解行列の和が階数1以下なら一方の因子が比例することから従う。
従って $w(q)$ にも一方の支持が高々1となる。
位置0では $L$ の全点で逆支持3、位置3では中央支持が少なくとも2なので矛盾する。

最後に $d=3$ では三列 $K_j$ は全点で独立である。これらが張る $3\mathcal O(-1)$ の
$4\mathcal O$ における商は $\mathcal O(3)$ で、切断を取れば
$H^0(4\mathcal O)\to H^0(\mathcal O(3))$ は同型である。
従って $y(t)$ は全位置2を張る三次曲線となる。一般の $s$ で $w(x(s))$ の
位置0,3の支持が次元3,4なので、$A(t)$ の左・右支持もそれぞれ少なくとも3,4である。
階数1の二項多項式行列を原始な一列と多項式行に因数分解すると、
左右支持の次元はそれぞれの因子の次数に1を加えた値以下である。
次数2では左因子に次数2以上、右因子に次数3以上を割り振れず矛盾する。
\end{proof}

\begin{corollary}\label{cor:rank-eight-center-types}
反対階数8で曲面射影がなければ、順・逆中心の中央支持は2または4である。
中央支持4の順中心の後続直線は、一つ前の位置の中央支持2の順中心の中央支持の消去直線に等しい。
中央支持2の順中心がある位置に他の順中心はない。特に一位置の順中心は高々一つである。
\end{corollary}
\begin{proof}
$a\ge3$ なら定理により $h=3$ であり、定理\ref{thm:branch-middle-defect-bound}の
反対点 $y$ は一つ前の位置で中央支持2である。$y$ が逆中心ならその対称二乗平面は
$R_1\cap(S\otimes V_2)$ 全体となり、$p$ の中央支持も2になって矛盾する。
ここで $S$ は $p$ の順支持である。$y$ は従って順中心で、その中央支持は $S$。
$R_0$ には $S$ の全方向に終わる純関係がある。$a=3$ なら $p$ の末尾からの
三次元縮約像は $2\times3$ 行列の中でSegreと交わり、連続純関係ができる。
従って $a=4$ である。

中央支持2の順中心を位置 $i$ に置く。その二項対称二乗の三関係により、辺 $i+1\to i+2$ では
第一点が後続直線上にあるか、第二点が中央支持の消去直線上にある。
同位置の別の順中心の後続直線は最初の直線と交わらないので、その第二曲線は直線になり、
次数式から二直線列となる。これは定理に反する。中央支持4の順中心が二つあれば、
それぞれの前位置の中央支持2の順中心は従って同一となり、消去直線も一致してしまう。
二つの順中心の後続直線が交わらないことに反する。
\end{proof}
